\scanpage{26}
\mspage{24.}
\uline{en schémas de groupes} \inserted{\uline{lisses}} $\underline{G}(A/S)$, et pour chaque prolongement \inserted{infinitésimal} $B$ de $A$ à un $T$, un isomorphisme canonique $\underline{G}(A/S)(T)\simeq G(B/T)$; tout ça bien sûr fonctoriel en $A$ et compatible avec changements de base. D'autre part, on peut préciser alors 3º en indiquant quel est le foncteur quasi-inverse de celui envisagé dans cet énoncé: l'extension de la filtration de $\operatorname{DR}^1(A/\unsure{S|A})(S)$ en une filtration du faisceau de Lie $\operatorname{DR}^1(A/S)(T)^{\vee}$ de $\underline{G}(A/S)(T)$ revient à étendre le sous-groupe vectoriel canonique de $\underline{G}(A/S)(S)$ en un sous-groupe vectoriel de $\underline{G}(A/S)(T)$, et l\deleted{a}'on trouve $B$ en passant simplement au quotient.

NB Je suis tombé sur la connexion canonique de $G(A/S)$ en essayant de simplifier la construction de Manin de l'application $A(\unsure{S|A})\longrightarrow\Gamma(S,\underline{O}_S)$ associée à une équation de Picart-Fuchs sur $S$, relativement à $A$. Pour ceci, il suffit de noter que localement sur $S$ toute section de $A$ sur $S$ se remonte en une section de $G=G(A/S)$ sur $S$? La donnée de la connexion de $G$ permet alors de prendre la dérivée de cette section, qui est un élément de $\Gamma(S,\underline{t}_G\otimes\underline{\Omega}^1_S)$, déterminé modulo une section de l'image du faisceau $\underline{t}^{\vee}_{A^*}$, correspondant à l'indétermination du relèvement d'\uline{u}ne section de $A$ en une section de $G$. Les équations de Picar\deleted{y}\inserted{t}-Fuchs sont définies tout juste pour arriver, d'une telle section de \deleted{\unsure{$\Gamma($|$\Gamma(S,$}} $\underline{t}_G\otimes\underline{\Omega}^1_S$ (qui en fait est un \inserted{co}\uline{cycle}, compte tenu de la connexion canonique de $\underline{t}_G$ provenant de $G$), avec l'indétermination qu'on vient de préciser, à tirer une section de $\underline{O}_S$ ... (Du moins en caractéristique nulle, cas dans lequel se place Manin de toutes façons).
\margin{25}
